% Generated draft: requires mathematical and source review.
\documentclass[11pt]{article}
\usepackage[margin=25mm]{geometry}
\usepackage{fontspec}
\setmainfont{Latin Modern Roman}
\usepackage{amsmath,amssymb,mathtools,unicode-math}
\setmathfont{Latin Modern Math}
\usepackage{enumitem,needspace,xurl,hyperref}
\setlength{\parindent}{0pt}
\setlength{\parskip}{5pt}
\setlength{\emergencystretch}{3em}
\newcommand{\N}{\mathbb{N}}
\newcommand{\Z}{\mathbb{Z}}
\newcommand{\Q}{\mathbb{Q}}
\newcommand{\R}{\mathbb{R}}
\newcommand{\C}{\mathbb{C}}
\newcommand{\F}{\mathbb{F}}
\newcommand{\A}{\mathbb{A}}
\newcommand{\PP}{\mathbb{P}}
\newcommand{\E}{\mathbb{E}}
\newcommand{\eps}{\varepsilon}
\newcommand{\dd}{\,\mathrm{d}}
\begin{document}
\section*{UnsolvedMath ID: 800275}
\textbf{Source problem code:} catalogue entry
\subsection{Definitions and mathematical statement}
Let \(\PP^2(\C)\) denote the complex projective plane. Let \(f:\PP^2(\C)\to\PP^2(\C)\) be a holomorphic map of degree \(d\geq 2\). The supplied source refers to an ergodic measure \(\mu=T\wedge T\), but does not define \(T\) or its normalization. Here \(T\wedge T\) denotes the wedge product of the intended current with itself; its definition and the hypotheses ensuring that it is an \(f\)-invariant ergodic probability measure require clarification.

For the intended measure \(\mu\), let \(\lambda_1\leq\lambda_2\) be the two complex Lyapunov exponents of the derivative cocycle of \(f\). Thus, when these exponents exist, if \(s_1(Df_x^n)\leq s_2(Df_x^n)\) are the singular values with respect to a smooth Hermitian metric, then for \(\mu\)-almost every \(x\),
\[
\lambda_i=\lim_{n\to\infty}\frac{1}{n}\log s_i(Df_x^n),\qquad i=1,2.
\]
The source question is whether, for every such \(f\) and its intended measure \(\mu=T\wedge T\),
\[
\lambda_1>0.
\]
This formulation remains conditional on identifying the intended current \(T\) and the meaning of degree. The supplied excerpt does not specify whether degree means algebraic degree or topological degree; these conventions differ for endomorphisms of the complex projective plane.
\subsection{Short English statement}
For a holomorphic self-map of the complex projective plane of degree at least two, must the smaller Lyapunov exponent of the intended ergodic measure \(T\wedge T\) be strictly positive? The supplied excerpt leaves \(T\) and the degree convention unspecified.
\subsection{Sources}
\begin{itemize}
\item Local UnsolvedMath registry: permanent ID \(800275\), exact source code \texttt{catalogue entry}; title: Several Complex Variables Problem List --- Problem 74. The registry attributes the problem to Buzzard, Fornaess, and Krantz, with proposed year 1995.
\item Buzzard, Fornaess, and Krantz, \emph{The Several Complex Variables Problem List} (1995), Problem 74. Supplied source URL: \url{https://arxiv.org/abs/math/9509202}. The supplied provenance records access on 2026-10-16 and extraction by \texttt{arxiv\_latex}; these are supplied provenance claims, not independently verified access.
\item The supplied registry labels this problem \texttt{open}. This is a source status label, not an independently verified assessment of the literature.
\end{itemize}
\end{document}
